% Versión simplificada del esquema del lazo de cálculo de la potencia límite por DNBR,
% pensada para una diapositiva de beamer.
% Requiere en el preámbulo:
%   \usepackage{tikz}
%   \usetikzlibrary{arrows.meta, positioning, shapes.geometric}
% Incluir en el frame con, por ejemplo:
%   \begin{frame}{Lazo de cálculo de la potencia límite por DNBR}
%       \centering
%       % Versión simplificada del esquema del lazo de cálculo de la potencia límite por DNBR,
% pensada para una diapositiva de beamer.
% Requiere en el preámbulo:
%   \usepackage{tikz}
%   \usetikzlibrary{arrows.meta, positioning, shapes.geometric}
% Incluir en el frame con, por ejemplo:
%   \begin{frame}{Lazo de cálculo de la potencia límite por DNBR}
%       \centering
%       % Versión simplificada del esquema del lazo de cálculo de la potencia límite por DNBR,
% pensada para una diapositiva de beamer.
% Requiere en el preámbulo:
%   \usepackage{tikz}
%   \usetikzlibrary{arrows.meta, positioning, shapes.geometric}
% Incluir en el frame con, por ejemplo:
%   \begin{frame}{Lazo de cálculo de la potencia límite por DNBR}
%       \centering
%       % Versión simplificada del esquema del lazo de cálculo de la potencia límite por DNBR,
% pensada para una diapositiva de beamer.
% Requiere en el preámbulo:
%   \usepackage{tikz}
%   \usetikzlibrary{arrows.meta, positioning, shapes.geometric}
% Incluir en el frame con, por ejemplo:
%   \begin{frame}{Lazo de cálculo de la potencia límite por DNBR}
%       \centering
%       \input{Figures/loop_potencia_dnbr_beamer.tex}
%   \end{frame}
\begin{tikzpicture}[
    font=\footnotesize,
    arr/.style={-{Latex[length=1.8mm]}, thick},
    node distance=6mm and 12mm,
    box/.style={
        rectangle, rounded corners, draw, thick,
        fill=blue!8, align=center,
        text width=34mm, minimum height=7mm, inner sep=1mm
    },
    startbox/.style={box, fill=gray!15},
    decision/.style={
        diamond, draw, thick, fill=orange!12, aspect=2.4,
        align=center, text width=22mm, inner sep=0mm
    },
    endbox/.style={
        rectangle, rounded corners, draw, thick,
        fill=green!12, align=center,
        text width=14mm, minimum height=7mm, inner sep=1mm
    },
]

    % Flujo principal
    \node[startbox]                  (init)   {Potencia semilla};
    \node[box, below=of init]        (caudal) {Caudal másico};
    \node[box, below=of caudal]      (cobra)  {COBRA3-CP};
    \node[box, below=of cobra]       (extrae) {Potencia a DNBR límite};
    \node[decision, below=of extrae] (conv)   {¿Converge?};
    \node[endbox, right=16mm of conv] (fin)   {Fin};

    \draw[arr] (init)   -- (caudal);
    \draw[arr] (caudal) -- (cobra);
    \draw[arr] (cobra)  -- (extrae);
    \draw[arr] (extrae) -- (conv);
    \draw[arr] (conv)   -- node[above] {Sí} (fin);

    % Realimentación: si no converge, se repite el cálculo de caudal
    \coordinate (loop) at ([xshift=-10mm] caudal.west);
    \draw[arr] (conv.west) -- node[above, pos=0.45] {No} (loop |- conv.west)
        -- (loop) -- (caudal.west);

\end{tikzpicture}

%   \end{frame}
\begin{tikzpicture}[
    font=\footnotesize,
    arr/.style={-{Latex[length=1.8mm]}, thick},
    node distance=6mm and 12mm,
    box/.style={
        rectangle, rounded corners, draw, thick,
        fill=blue!8, align=center,
        text width=34mm, minimum height=7mm, inner sep=1mm
    },
    startbox/.style={box, fill=gray!15},
    decision/.style={
        diamond, draw, thick, fill=orange!12, aspect=2.4,
        align=center, text width=22mm, inner sep=0mm
    },
    endbox/.style={
        rectangle, rounded corners, draw, thick,
        fill=green!12, align=center,
        text width=14mm, minimum height=7mm, inner sep=1mm
    },
]

    % Flujo principal
    \node[startbox]                  (init)   {Potencia semilla};
    \node[box, below=of init]        (caudal) {Caudal másico};
    \node[box, below=of caudal]      (cobra)  {COBRA3-CP};
    \node[box, below=of cobra]       (extrae) {Potencia a DNBR límite};
    \node[decision, below=of extrae] (conv)   {¿Converge?};
    \node[endbox, right=16mm of conv] (fin)   {Fin};

    \draw[arr] (init)   -- (caudal);
    \draw[arr] (caudal) -- (cobra);
    \draw[arr] (cobra)  -- (extrae);
    \draw[arr] (extrae) -- (conv);
    \draw[arr] (conv)   -- node[above] {Sí} (fin);

    % Realimentación: si no converge, se repite el cálculo de caudal
    \coordinate (loop) at ([xshift=-10mm] caudal.west);
    \draw[arr] (conv.west) -- node[above, pos=0.45] {No} (loop |- conv.west)
        -- (loop) -- (caudal.west);

\end{tikzpicture}

%   \end{frame}
\begin{tikzpicture}[
    font=\footnotesize,
    arr/.style={-{Latex[length=1.8mm]}, thick},
    node distance=6mm and 12mm,
    box/.style={
        rectangle, rounded corners, draw, thick,
        fill=blue!8, align=center,
        text width=34mm, minimum height=7mm, inner sep=1mm
    },
    startbox/.style={box, fill=gray!15},
    decision/.style={
        diamond, draw, thick, fill=orange!12, aspect=2.4,
        align=center, text width=22mm, inner sep=0mm
    },
    endbox/.style={
        rectangle, rounded corners, draw, thick,
        fill=green!12, align=center,
        text width=14mm, minimum height=7mm, inner sep=1mm
    },
]

    % Flujo principal
    \node[startbox]                  (init)   {Potencia semilla};
    \node[box, below=of init]        (caudal) {Caudal másico};
    \node[box, below=of caudal]      (cobra)  {COBRA3-CP};
    \node[box, below=of cobra]       (extrae) {Potencia a DNBR límite};
    \node[decision, below=of extrae] (conv)   {¿Converge?};
    \node[endbox, right=16mm of conv] (fin)   {Fin};

    \draw[arr] (init)   -- (caudal);
    \draw[arr] (caudal) -- (cobra);
    \draw[arr] (cobra)  -- (extrae);
    \draw[arr] (extrae) -- (conv);
    \draw[arr] (conv)   -- node[above] {Sí} (fin);

    % Realimentación: si no converge, se repite el cálculo de caudal
    \coordinate (loop) at ([xshift=-10mm] caudal.west);
    \draw[arr] (conv.west) -- node[above, pos=0.45] {No} (loop |- conv.west)
        -- (loop) -- (caudal.west);

\end{tikzpicture}

%   \end{frame}
\begin{tikzpicture}[
    font=\footnotesize,
    arr/.style={-{Latex[length=1.8mm]}, thick},
    node distance=6mm and 12mm,
    box/.style={
        rectangle, rounded corners, draw, thick,
        fill=blue!8, align=center,
        text width=34mm, minimum height=7mm, inner sep=1mm
    },
    startbox/.style={box, fill=gray!15},
    decision/.style={
        diamond, draw, thick, fill=orange!12, aspect=2.4,
        align=center, text width=22mm, inner sep=0mm
    },
    endbox/.style={
        rectangle, rounded corners, draw, thick,
        fill=green!12, align=center,
        text width=14mm, minimum height=7mm, inner sep=1mm
    },
]

    % Flujo principal
    \node[startbox]                  (init)   {Potencia semilla};
    \node[box, below=of init]        (caudal) {Caudal másico};
    \node[box, below=of caudal]      (cobra)  {COBRA3-CP};
    \node[box, below=of cobra]       (extrae) {Potencia a DNBR límite};
    \node[decision, below=of extrae] (conv)   {¿Converge?};
    \node[endbox, right=16mm of conv] (fin)   {Fin};

    \draw[arr] (init)   -- (caudal);
    \draw[arr] (caudal) -- (cobra);
    \draw[arr] (cobra)  -- (extrae);
    \draw[arr] (extrae) -- (conv);
    \draw[arr] (conv)   -- node[above] {Sí} (fin);

    % Realimentación: si no converge, se repite el cálculo de caudal
    \coordinate (loop) at ([xshift=-10mm] caudal.west);
    \draw[arr] (conv.west) -- node[above, pos=0.45] {No} (loop |- conv.west)
        -- (loop) -- (caudal.west);

\end{tikzpicture}
